%==================== 数学符号公式 ============
\usepackage{amsmath}
\usepackage[style=1]{mdframed}
\usepackage{amsthm}
\usepackage{pdfpages}
\usepackage{amsfonts}
\usepackage{mathrsfs}
\usepackage{bm}
\usepackage{manfnt}
\usepackage{lettrine}
\def\attention{\lettrine[lines=2,lraise=0,nindent=0em]{\large\textdbend\hspace{1mm}}{}}
\usepackage{longtable}
\usepackage[toc,page]{appendix}
\usepackage{geometry}
\geometry{top=3.0cm,bottom=2.7cm,left=2.5cm,right=2.5cm} 

\usepackage{fancyhdr}
%====================清单==========================
\usepackage{enumitem}
\setlist[description]{leftmargin=\parindent,labelindent=\parindent}
%====================公式按章编号==========================
\numberwithin{equation}{section}
\numberwithin{table}{section}
\numberwithin{figure}{section}
%================= 基本格式预置 ===========================

\pagestyle{fancy}
\fancyhf{}
\fancyfoot[C]{~\zihao{5} \thepage~}
\renewcommand{\headrulewidth}{0.65pt}

% 修正：使用 \ctexset 替代已弃用的 \CTEXsetup
\ctexset{
  section = {
    format = {\centering\bfseries\zihao{3}},
    name = {第, 节},
    number = \arabic{section},
  },
  subsection = {
    format = {\bfseries\zihao{4}},
    name = {},
    number = \arabic{subsection},
  },
  subsubsection = {
    format = {\songti\zihao{-4}},
    name = {},
    number = \arabic{subsubsection},
  }
}

%================== 表格 =========================
\usepackage{makecell,rotating,multirow,diagbox}
\newcommand{\addtbl}[1]
{
	\input{table/#1}
}
%==================  图片 =========================
\newcommand{\addfig}[4][0.5]
{
  \begin{figure}[thbp!]
  \centering\includegraphics[width=#1\linewidth]{figure/#2}
  \caption{#4}
  \label{#3}
  \end{figure}
}
%================== 图形支持宏包 =========================
\usepackage{subfigure}
\usepackage{graphicx}
\usepackage{color,xcolor}
\usepackage{hyperref}
\usepackage[font=small,labelfont=bf,labelsep=space]{caption}
\captionsetup{figurewithin=section}
%==================== 源码和流程图 =====================
\newcommand{\addcode}[2][None]
{
	\lstinputlisting[language=#1]{code/#2}
}

\usepackage{listings}
\usepackage{xcolor}
\usepackage{color}
\definecolor{dkgreen}{rgb}{0,0.6,0}
\definecolor{gray}{rgb}{0.5,0.5,0.5}
\definecolor{mauve}{rgb}{0.58,0,0.82}
\definecolor{comment}{rgb}{0.56,0.64,0.68}

\lstset{
  frame=tb,
  aboveskip=3mm,
  belowskip=3mm,
  xleftmargin=2em,
  xrightmargin=2em,
  showstringspaces=false,
  keepspaces=true,
  columns=flexible,
  framerule=1pt,
  rulecolor=\color{gray!35},
  backgroundcolor=\color{gray!5},
  basicstyle={\small\ttfamily},
  numbers=none,
  numberstyle=\tiny\color{gray},
  keywordstyle=\color{blue},
  commentstyle=\color{comment},
  stringstyle=\color{dkgreen},
  breaklines=true,
  breakatwhitespace=true,
  tabsize=2,
}
%--------------------
\hypersetup{hidelinks}
\usepackage{booktabs}  
\usepackage{shorttoc}
\usepackage{tabu,tikz}
\usepackage{float}
\usepackage{multirow}

\tabcolsep=1ex
\tabulinesep=\tabcolsep
\newlength\tikzboxwidth
\newlength\tikzboxheight
\newcommand\tikzbox[1]{%
        \settowidth\tikzboxwidth{#1}%
        \settoheight\tikzboxheight{#1}%
        \begin{tikzpicture}
        \path[use as bounding box]
                (-0.5\tikzboxwidth,-0.5\tikzboxheight)rectangle
                (0.5\tikzboxwidth,0.5\tikzboxheight);
        \node[inner sep=\tabcolsep+0.5\arrayrulewidth,line width=0.5mm,draw=black]
                at(0,0){#1};
        \end{tikzpicture}%
        }

\makeatletter
\def\hlinew#1{%
  \noalign{\ifnum0=`}\fi\hrule \@height #1 \futurelet
   \reserved@a\@xhline}
   
\newcommand{\tabincell}[2]{\begin{tabular}{@{}#1@{}}#2\end{tabular}}%

\usepackage{subfigure}
\usepackage{CJK}
\usepackage{ifthen}

\newcommand{\HRule}{\rule{\linewidth}{0.5mm}}

\newtheorem{Theorem}{定理}
\newtheorem{Lemma}{引理}

\makeatletter
\@addtoreset{equation}{section}
\makeatother
\renewcommand{\theequation}{\arabic{section}.\arabic{equation}}

%-------------------------
\usepackage{pythonhighlight}
\usepackage{tikz}                    
\usepackage{tikz-3dplot}
\usetikzlibrary{shapes,arrows,positioning}

% 参考文献格式
\bibliographystyle{gbt7714-2005}

%===================  定理类环境定义 ===================
\newtheorem{example}{例}
\newtheorem{algorithm}{算法}
\newtheorem{theorem}{定理}
\newtheorem{definition}{定义}
\newtheorem{axiom}{公理}
\newtheorem{property}{性质}
\newtheorem{proposition}{命题}
\newtheorem{lemma}{引理}
\newtheorem{corollary}{推论}
\newtheorem{remark}{注解}
\newtheorem{condition}{条件}
\newtheorem{conclusion}{结论}
\newtheorem{assumption}{假设}

%==================重定义 ===================
\renewcommand{\contentsname}{目录}     
\renewcommand{\abstractname}{摘要}   
\renewcommand{\indexname}{索引}
\renewcommand{\figurename}{图}
\renewcommand{\tablename}{表}
\renewcommand{\appendixname}{附录}
\renewcommand{\proofname}{证明}